
\documentclass[12pt,a4paper]{article}

% ------------------------------------------------
% Paquetes
% ------------------------------------------------
\usepackage[utf8]{inputenc}
\usepackage[T1]{fontenc}
\usepackage[spanish, es-nodecimaldot]{babel}
\usepackage{geometry}
\usepackage{graphicx}
\usepackage{float}
\usepackage{xcolor}
\usepackage{listings}
\usepackage{hyperref}
\usepackage{fancyhdr}
\usepackage{titlesec}

% ------------------------------------------------
% Configuración de página
% ------------------------------------------------
\geometry{
    top=2.5cm,
    bottom=2.5cm,
    left=2.5cm,
    right=2.5cm
}

% ------------------------------------------------
% Configuración de listings
% ------------------------------------------------
\definecolor{codegray}{gray}{0.95}

\lstset{
    backgroundcolor=\color{codegray},
    basicstyle=\ttfamily\small,
    breaklines=true,
    frame=single,
    columns=fullflexible,
    keepspaces=true,
    showstringspaces=false,
    tabsize=2
}

% ------------------------------------------------
% Configuración de títulos
% ------------------------------------------------
\titleformat{\section}
    {\large\bfseries}
    {\thesection.}
    {0.5em}
    {}

% ------------------------------------------------
% Hipervínculos
% ------------------------------------------------
\hypersetup{
    colorlinks=true,
    linkcolor=black,
    urlcolor=blue,
    citecolor=black
}

% ------------------------------------------------
% Comando para insertar capturas opcionales
% ------------------------------------------------
\newcommand{\captura}[2]{%
    \begin{figure}[H]
        \centering
        \IfFileExists{#1}{%
            \includegraphics[
                width=0.92\textwidth,
                keepaspectratio
            ]{#1}
        }{%
            \fbox{%
                \parbox[c][4cm][c]{0.82\textwidth}{%
                    \centering
                    \textbf{Captura pendiente}\\[0.5em]
                    Archivo esperado: \texttt{#1}
                }%
            }%
        }
        \caption{#2}
    \end{figure}
}

% ------------------------------------------------
% Datos del documento
% ------------------------------------------------
\title{
    MPI:\\[0.3em]
    \textbf{Suma de elementos de un arreglo}\\[0.5em]
    \large Supercómputo
}

\author{
    Jorge Alberto Suárez Saldaña\\
    \small Matrícula: 355992\\
    \small Profesor: Raymundo Antonio González Grimaldo
}

\date{\today}

% ------------------------------------------------
% Documento
% ------------------------------------------------
\begin{document}

\maketitle

\section{Introducción}
En la presente tarea se utilizó MPI en el lenguaje de programación C para paralelizar el algoritmo para la suma de los elementos de un arreglo dividiendo la información en subarreglos que se le pasaban a subprocesos, utilizando el proceso cero como proceso controlador para enviar la información a los procesos esclavos y calcular la suma de las sumas parciales obtenidas ñor los procesos esclavos.

\section{Código fuente}
\begin{verbatim}
#include <stdio.h>
#include <stdlib.h>
#include <time.h>
#include <mpi.h>

#define N 200

int main(int argc, char *argv[]) {
    int idProc, numProc;
    int nDatos, nDatosU;
    int suma = 0, sumaP = 0;
    int *a;

    MPI_Init(&argc, &argv);

    MPI_Comm_rank(MPI_COMM_WORLD, &idProc);
    MPI_Comm_size(MPI_COMM_WORLD, &numProc);

    if (numProc < 2 || numProc > N + 1) {
        if (idProc == 0)
            printf("Se requieren entre 2 y %d procesos.\n", N + 1);

        MPI_Finalize();
        return 1;
    }

    nDatos = N / (numProc - 1);
    nDatosU = nDatos + N % (numProc - 1);

    if (idProc == 0) {
        // Proceso coordinador
        a = malloc(N * sizeof(int));

        if (a == NULL) {
            MPI_Abort(MPI_COMM_WORLD, 1);
        }

        // Llenar arreglo con numeros aleatorios del 1 al 100
        srand(time(NULL));

        for (int i = 0; i < N; i++) {
            a[i] = rand() % 100 + 1;
        }

        // Imprimir arreglo original
        printf("Arreglo original:\n");

        for (int i = 0; i < N; i++) {
            printf("%d ", a[i]);
        }

        printf("\n\n");

        // Distribuir datos a los trabajadores
        for (int iProc = 1; iProc < numProc - 1; iProc++) {
            MPI_Send(
                &a[(iProc - 1) * nDatos],
                nDatos,
                MPI_INT,
                iProc,
                0,
                MPI_COMM_WORLD
            );
        }

        // Enviar la ultima seccion
        MPI_Send(
            &a[(numProc - 2) * nDatos],
            nDatosU,
            MPI_INT,
            numProc - 1,
            0,
            MPI_COMM_WORLD
        );

        // Recibir sumas parciales
        for (int iProc = 1; iProc < numProc; iProc++) {
            MPI_Recv(
                &sumaP,
                1,
                MPI_INT,
                iProc,
                0,
                MPI_COMM_WORLD,
                MPI_STATUS_IGNORE
            );

            printf("Proceso %d: suma parcial = %d\n",
                   iProc, sumaP);

            suma += sumaP;
        }

        printf("\nSuma total paralela: %d\n", suma);

        // Verificacion secuencial
        int sumaSecuencial = 0;

        for (int i = 0; i < N; i++) {
            sumaSecuencial += a[i];
        }

        printf("Suma secuencial: %d\n", sumaSecuencial);

        if (suma == sumaSecuencial)
            printf("Resultado correcto.\n");
        else
            printf("Error: las sumas no coinciden.\n");

        free(a);

    } else {
        // Procesos trabajadores
        if (idProc == numProc - 1) {
            nDatos = nDatosU;
        }

        a = malloc(nDatos * sizeof(int));

        if (a == NULL) {
            MPI_Abort(MPI_COMM_WORLD, 1);
        }

        // Recibir seccion del arreglo
        MPI_Recv(
            a,
            nDatos,
            MPI_INT,
            0,
            0,
            MPI_COMM_WORLD,
            MPI_STATUS_IGNORE
        );

        // Calcular suma parcial
        for (int i = 0; i < nDatos; i++) {
            suma += a[i];
        }

        // Enviar resultado al coordinador
        MPI_Send(
            &suma,
            1,
            MPI_INT,
            0,
            0,
            MPI_COMM_WORLD
        );

        free(a);
    }

    MPI_Finalize();
    return 0;
}

\end{verbatim}

\section{Explicación del código}

El programa comienza incluyendo las bibliotecas necesarias para las operaciones de entrada y salida, administración de memoria, generación de números aleatorios y comunicación entre procesos.

La constante \texttt{N} establece que el arreglo tendrá 200 elementos.

\begin{lstlisting}[language=C]
#define N 200

MPI_Init(&argc, &argv);

MPI_Comm_rank(MPI_COMM_WORLD, &idProc);
MPI_Comm_size(MPI_COMM_WORLD, &numProc);
\end{lstlisting}

La función \texttt{MPI\_Init()} inicializa el entorno de ejecución MPI.

Posteriormente, \texttt{MPI\_Comm\_rank()} obtiene el identificador del proceso actual y lo almacena en \texttt{idProc}, mientras que \texttt{MPI\_Comm\_size()} obtiene el número total de procesos participantes y lo guarda en \texttt{numProc}.

El comunicador \texttt{MPI\_COMM\_WORLD} representa al conjunto inicial de procesos que participan en la ejecución.

El programa comprueba que existan al menos dos procesos y que no se supere el límite de 201. Esto permite disponer de un coordinador y, como máximo, un trabajador por elemento del arreglo.

\subsection{Distribución de los datos}

Para repartir el trabajo, el programa calcula la cantidad de elementos que recibirá cada proceso trabajador:

\begin{lstlisting}[language=C]
nDatos = N / (numProc - 1);
nDatosU = nDatos + N % (numProc - 1);
\end{lstlisting}

El proceso cero no participa directamente en la suma de las secciones del arreglo, por lo que el número de trabajadores es:

\[
P_t = P - 1
\]

donde \(P\) representa el número total de procesos.

La cantidad de elementos asignada a cada trabajador, excepto al último, se obtiene mediante división entera:

\[
nDatos = \left\lfloor \frac{N}{P-1} \right\rfloor
\]

Cuando el número de elementos no es divisible exactamente entre la cantidad de trabajadores, existe un residuo. Para evitar que algunos elementos queden sin procesar, dicho residuo se agrega a la sección del último trabajador:

\[
nDatosU =
\left\lfloor \frac{N}{P-1} \right\rfloor
+
\left(N \bmod (P-1)\right)
\]

Por ejemplo, si se utilizan siete procesos, seis de ellos serán trabajadores:

\[
nDatos = \left\lfloor \frac{200}{6} \right\rfloor = 33
\]

\[
nDatosU = 33 + (200 \bmod 6) = 35
\]

De esta manera, cinco trabajadores procesarán 33 elementos cada uno y el último procesará 35, cubriendo los 200 elementos originales.

\subsection{Funcionamiento del proceso coordinador}

El proceso con identificador cero es responsable de administrar los datos y reunir los resultados.

Inicialmente reserva memoria dinámica mediante \texttt{malloc()} para almacenar los 200 enteros del arreglo. Si ocurre un error durante la asignación de memoria, se utiliza \texttt{MPI\_Abort()} para interrumpir la ejecución MPI.

Posteriormente, el arreglo se llena con números pseudoaleatorios comprendidos entre 1 y 100:

\begin{lstlisting}[language=C]
srand(time(NULL));

for (int i = 0; i < N; i++) {
    a[i] = rand() % 100 + 1;
}
\end{lstlisting}

La función \texttt{srand()} inicializa la semilla del generador pseudoaleatorio utilizando la hora actual. La expresión \texttt{rand() \% 100 + 1} produce valores enteros dentro del intervalo indicado.

Una vez generado el arreglo, el coordinador imprime su contenido y comienza a distribuir los datos.

\subsubsection{Envío de las secciones}

La distribución se realiza mediante la función \texttt{MPI\_Send()}:

\begin{lstlisting}[language=C]
for (int iProc = 1; iProc < numProc - 1; iProc++) {
    MPI_Send(
        &a[(iProc - 1) * nDatos],
        nDatos,
        MPI_INT,
        iProc,
        0,
        MPI_COMM_WORLD
    );
}
\end{lstlisting}

Cada trabajador recibe una sección consecutiva del arreglo. La expresión

\[
(iProc-1)\cdot nDatos
\]

determina el índice inicial de la sección correspondiente a cada proceso.

Los parámetros principales de \texttt{MPI\_Send()} indican la dirección inicial de los datos, la cantidad de elementos, su tipo, el proceso destinatario, la etiqueta del mensaje y el comunicador.

Para el último trabajador se realiza un envío independiente, utilizando \texttt{nDatosU} como cantidad de elementos. Esto garantiza que también se procesen los elementos sobrantes de la división.

\subsubsection{Recepción de las sumas parciales}

Después de distribuir el arreglo, el coordinador recibe los resultados calculados por los trabajadores:

\begin{lstlisting}[language=C]
for (int iProc = 1; iProc < numProc; iProc++) {
    MPI_Recv(
        &sumaP,
        1,
        MPI_INT,
        iProc,
        0,
        MPI_COMM_WORLD,
        MPI_STATUS_IGNORE
    );

    suma += sumaP;
}
\end{lstlisting}

La función \texttt{MPI\_Recv()} espera y recibe un mensaje procedente del trabajador indicado.

En este caso, cada mensaje contiene un único entero correspondiente a la suma parcial. Dicho resultado se almacena temporalmente en \texttt{sumaP} y posteriormente se acumula en \texttt{suma}.

La constante \texttt{MPI\_STATUS\_IGNORE} indica que no es necesario conservar información adicional sobre el mensaje recibido.

El resultado final se expresa como:

\[
S = \sum_{j=1}^{P-1} S_j
\]

donde \(S_j\) representa la suma parcial calculada por el trabajador \(j\).

\subsection{Funcionamiento de los procesos trabajadores}

Todos los procesos cuyo identificador es diferente de cero ejecutan la sección del programa destinada a los trabajadores.

Primero, cada trabajador determina la cantidad de elementos que debe procesar. Si se trata del último proceso, utiliza \texttt{nDatosU}; de lo contrario, conserva el valor de \texttt{nDatos}.

Después reserva memoria para su arreglo local y recibe la sección enviada por el coordinador:

\begin{lstlisting}[language=C]
MPI_Recv(
    a,
    nDatos,
    MPI_INT,
    0,
    0,
    MPI_COMM_WORLD,
    MPI_STATUS_IGNORE
);
\end{lstlisting}

Es importante señalar que cada proceso cuenta con su propio espacio de memoria. Por esta razón, el arreglo local del trabajador es independiente del arreglo original almacenado por el coordinador.

Una vez recibidos los elementos, el trabajador realiza una suma secuencial sobre su sección:

\begin{lstlisting}[language=C]
for (int i = 0; i < nDatos; i++) {
    suma += a[i];
}
\end{lstlisting}

Aunque esta operación utiliza un ciclo secuencial, los diferentes trabajadores pueden ejecutar sus cálculos de forma concurrente.

Finalmente, cada trabajador envía su suma parcial al proceso cero:

\begin{lstlisting}[language=C]
MPI_Send(
    &suma,
    1,
    MPI_INT,
    0,
    0,
    MPI_COMM_WORLD
);
\end{lstlisting}

El trabajador libera posteriormente la memoria dinámica mediante \texttt{free()}.

\subsection{Verificación de resultados}

Para comprobar la exactitud del cálculo paralelo, el proceso coordinador realiza una segunda suma utilizando un algoritmo secuencial:

\begin{lstlisting}[language=C]
int sumaSecuencial = 0;

for (int i = 0; i < N; i++) {
    sumaSecuencial += a[i];
}
\end{lstlisting}

Se comparan los resultados obtenidos por ambos métodos:

Si los valores coinciden, el programa imprime el mensaje \texttt{Resultado correcto.}; de lo contrario, informa que existe una diferencia entre las sumas.

Esta verificación permite identificar posibles errores en la distribución de los elementos, las comunicaciones o la acumulación de los resultados parciales.

\subsection{Finalización y liberación de recursos}

Al terminar sus operaciones, cada proceso libera la memoria dinámica que haya reservado.

Posteriormente, todos los procesos ejecutan:

\begin{lstlisting}[language=C]
MPI_Finalize();
\end{lstlisting}

Esta función finaliza ordenadamente el entorno MPI y permite liberar los recursos internos asociados a la ejecución paralela.

\subsection{Conclusión}

El programa implementa satisfactoriamente el esquema de paralelización coordinador-trabajadores mediante MPI. El proceso cero administra la distribución del arreglo y la recolección de resultados, mientras que los demás procesos calculan simultáneamente sumas parciales sobre diferentes secciones de los datos.

La implementación permite comprender el funcionamiento de las comunicaciones punto a punto mediante \texttt{MPI\_Send()} y \texttt{MPI\_Recv()}, así como la identificación y organización de procesos mediante \texttt{MPI\_Comm\_rank()} y \texttt{MPI\_Comm\_size()}.

La comparación con una implementación secuencial constituye un mecanismo de verificación que permite comprobar la corrección del resultado.

\section{Capturas de comprobación}
\captura{codigoFuente.png}{
    Codigo fuente en el editor NeoVim.
}
\captura{diferenteN.png}{
    Código fuente con una N diferente (200) correspondiente a segunda ejecución.
}
\captura{terminal.png}{
    Compilación y ejecución del programa con diferentes N y número de procesos.
}

\end{document}
